%% GENERATED by python/paper/build.py from results/ -- do not edit by hand.
%% Rebuild:  .venv\Scripts\python.exe python\paper\build.py --only US_ESC_Country
%% Source:   results/esc/escCountry.csv
\begin{table}[!htb]
\centering
\begin{threeparttable}
\caption{The UK and France under the US cost of redistribution}
\label{table:US_ESC:country}
\renewcommand{\arraystretch}{1.25}
\begin{tabularx}{\textwidth}{lYYYYY}
\toprule
\textbf{Economy} & $\theta^{\ast}$ \textbf{observed} & \textbf{Tax rate} & $\theta$ \textbf{chosen, US} $\lambda$ & \textbf{Own} $\lambda$ & $\tilde V$ \\
\midrule 
UK & 0.560 & 11.9\% & 0.625 & 7.264 & 0.223 \\
France & 1.000 & 21.3\% & 0.749 & -- & 0.218 \\
UK at US income groups & 0.543 & 11.9\% & 0.644 & 6.519 & 0.065 \\
\bottomrule
\end{tabularx}
\begin{tablenotes}[flushleft]
\footnotesize
\item[] \textit{Note:} The deadweight cost $f(\theta_t, \tau_t)$ of \eqref{eq:esc:budget} is quadratic in the implicit tax the flat component levies and scaled by the size of the system, $f = \exp\{-\tfrac12 \lambda \tau_t \tilde V (1-\theta_t)^2\}$, with $\lambda$ calibrated per $\rho$ (table \ref{table:US_ESC:calibration}). Each economy is its own calibration at $\rho = 1$ ($\beta$ imposed from the US, $\omega$ its own; table \ref{table:US:Calib}); the UK at US income groups is the UK workbook regrouped at the US income percentiles. ``Chosen'' is the design in force in 2020 on the economy\textquotesingle s own freely simulated path under the US parameter $\lambda = 8.643$, to be read against the observed one; ``own'' is the value at which that path re-elects the observed design, with $\omega$ recalibrated at each trial value. No finite $\lambda$ makes France re-elect its observed design: it is the corner $\theta = 1$, and under a cost quadratic in the redistribution performed the first unit of redistribution is free at the margin.
\end{tablenotes}
\end{threeparttable}
\end{table}
